\section{Related Work}
\label{sec:related}

\subsection{Short-Term Load Forecasting Methods and Benchmarks}
Short-term load forecasting (STLF) has a long history in power systems. Early reviews already questioned whether neural networks consistently outperform well-specified linear and exponential-smoothing models~\cite{hippert2001nn}, and empirical comparisons on European data showed that simple seasonal methods are hard to beat at short lead times~\cite{taylor2007european,taylor2010triple}. Subsequent surveys catalogue a wide range of statistical and machine-learning (ML) approaches~\cite{kuster2017review,haben2021lv,wang2019smartmeter,hong2020outlook}. The Global Energy Forecasting Competitions (GEFCom) popularised gradient-boosted trees~\cite{lloyd2014gefcom2012} and regularised linear models with rich calendar structure~\cite{ziel2016lasso} for load, and shifted attention to probabilistic forecasting~\cite{hong2016gefcom2014,hong2016tutorial,hong2019gefcom2017}. Gradient boosting, in particular LightGBM~\cite{ke2017lightgbm}, also dominated the M5 competition~\cite{makridakis2022m5}, and tree ensembles remain strong on tabular problems~\cite{grinsztajn2022tabular}.

A recurring criticism of the STLF literature is that new models are often evaluated on short, proprietary or idiosyncratically pre-processed data sets, against weak baselines, and without statistical testing~\cite{hong2020outlook,petropoulos2022forecasting}. In the related field of electricity price forecasting, Lago~\emph{et~al.}~\cite{lago2021benchmark} addressed this by releasing an open benchmark in which a well-regularised, hour-specific linear autoregressive model (LEAR) is shown to be competitive with deep networks, and by advocating rolling-window recalibration, long test periods and Diebold--Mariano (DM) testing~\cite{diebold1995dm,harvey1997hln,ziel2018multivariate}. We follow these best practices here for load.

\subsection{Calibration Windows and Concept Drift}
Under non-stationarity, the amount of history used to fit a model (the \emph{calibration window}) and how often the model is refitted are key design choices~\cite{gama2014drift}. For electricity prices, Marcjasz~\emph{et~al.}~\cite{marcjasz2018windows} showed that the choice of window has a substantial impact, and that averaging forecasts obtained from short and long windows is a robust alternative to selecting a single window~\cite{hubicka2019windows,serafin2019windows}. For load, the COVID-19 lockdowns triggered interest in adaptive methods that track abrupt changes~\cite{obst2021covid}. The rapid diffusion of behind-the-meter photovoltaic (PV) generation is a slower but persistent source of drift, because it continuously reshapes the daily profile of the load observed by the transmission system operator. To the best of our knowledge, no study has quantified the effect of calibration windows and refitting frequencies on day-ahead load forecasts in a system undergoing such a transition.

\subsection{Probabilistic Forecasting and Conformal Prediction}
Probabilistic forecasts are typically produced by quantile regression~\cite{koenker1978quantiles}, either with linear models~\cite{ziel2016lasso} or with gradient boosting, and evaluated with proper scoring rules such as the pinball loss and the Winkler interval score~\cite{gneiting2007scoring,winkler1972interval,nowotarski2018probabilistic}. Conformal prediction~\cite{vovk2022alrw} provides distribution-free coverage guarantees under exchangeability and can be combined with quantile regression~\cite{romano2019cqr}. Because time series are not exchangeable, adaptive conformal inference (ACI) updates the nominal miscoverage level online to retain long-run coverage under distribution shift~\cite{gibbs2021aci}, and has been applied to electricity data~\cite{zaffran2022aci}.

\subsection{Load Forecasting in Brazil}
In Brazil, the National System Operator (ONS) produces the official short-term load forecast that feeds the daily operation programming and the DESSEM hydrothermal scheduling model. Aquila~\emph{et~al.}~\cite{aquila2023brazil} review ML methods for STLF, including applications to the Brazilian system. The studies we are aware of typically address a single subsystem or a short test period, and do not compare their models with the forecasts actually used in system operation. Furthermore, the legal framework for distributed micro- and mini-generation (MMGD), consolidated by Law~14.300/2022~\cite{lei14300}, triggered an unprecedented expansion of rooftop PV, whose effect on the predictability of the operator's load has not been quantified. This paper fills these gaps with an open, reproducible benchmark covering all four subsystems of the National Interconnected System (SIN) and including the official ONS day-ahead forecast as a reference.
